\scanpage{97}
\handwritten{15.9.67.}\\

Cher Dieudonné,\\

Tes objections de ta lettre du 11 Septembre sont encore fondées. D'ailleurs, il ne me semble pas évident que le fait de pouvoir trouver $p,q$ fixes tels qu'on ait des suites exactes $A_n^p\quad A_n^q\quad M_n\quad0$, permette de s'en tirer; tant mieux si tu y arrives. L'existence de ces $p$ et $q$ me semble d'autre part facile, en utlisant le

\uline{Lemme}\quad Soit $A$ un anneau, $M$ un $A$-module, $(f_i)_{1\;i\;n}$ des éléments de $A$ engendrant l'idéal unité, \deleted{$m$ un entier} supposons que pour tout $i$, le $A_{f_i}$-module $M_{f_i}$ soit engendré par $m_i$ générateurs. Alors $M$ est engendré par $m\quad m_i$ générateurs.

Démonstration: soient $g_j^i\quad M_{f_i}\quad\unsure{(}1\quad j\quad m_i)$ les générateurs de $M_{f_i}$ , on aura donc $g_j^i=h_j^i/f_i^{N_{ij}}$ , avec $h_j^i\quad \unsure{\deleted{A}M|\deleted{M}A}$ . Alors les $\unsure{\deleted{g}h_j^i|\deleted{h}g_j^i}$ pour $j$ fixes engendrent $\overset{\handwritten{\sim}}{M}$ sur l'ouvert $\operatorname{Spec}(A_{f_i})$ de $\operatorname{Spec}(A_f)$, donc comme pour $i$ varaible ces ouverts recouvrent $\operatorname{Spec}(A)$, \unsure{\deleted{n}i|\deleted{o}i|\deleted{h}i|\deleted{i}n|\deleted{i}o|\deleted{i}h}l s'ensuit que les $h_j^i$ pour $i,j$ variables engendrent $\overset{\handwritten{\unsure{\sim}}}{M}$, donc $M$, ce qui établit le lemme.

Revenant alors à la situation de 10.10.5 , on sait que $M=\lim M_n$ est un $A$-module de type fini, considérons alors un ép\unsure{\deleted{m}i|\deleted{t}i|\deleted{i}m|\deleted{i}t}morphisme $u:A^q\quad M$ . Soient $f_i\quad A$ dont les images dans $A_o$ engendrent l'idéal unité, et tel que sur les \deleted{Spef(} ouverts correspondants $X_i$ de $X=\operatorname{Specf}(A)$, $F$ admette une présentation finie, donc soit de la forme $M_i$ , où $M_i$ est un $\inserted{A_i=}A_{f_i}$-\unsure{\deleted{ù}m|\deleted{à}m|\deleted{m}ù|\deleted{m}à}odul\unsure{e} de présentation finie. Alors $u$ restreint à $X_i$ définit un épimorphisme $A_i^q\quad M_i$ , et comme $M_i$ est de présentation finie, il existe un homomorphi\unsure{s}me $v_i$ rendant exacte la suite $A_i^p\quad A_i^q\quad M_i\quad0$ . Tensorisant par $A_o$ , on en déduit une suite exacte $(A_n)_{f_i}^p\quad(A_n)_{f_i}^q\quad(M_n)_{f_i}\quad0$, ce qui montre que,\deleted{le} si $R_n=\operatorname{Ker}(A_n^q\quad M_n)$, alors $R_n$ est sur $\operatorname{Spec}(A_{n\;f_i})$ engendré par $p$ éléments. Donc en vertu du lemme $R_n$ est engendré par $mp$ éléments
% archive sheet number: 96
